\subsection{معماری و پیاده‌سازی}
کلاس‌های اصلی این موتور شامل \lr{ToMoELinear} و \lr{ToMoECompressor} بوده و در رجیستری فشرده‌سازی سیستم تحت نام‌های \lr{tomoe} و \lr{tomoe\_compressor} به ثبت رسیده‌اند. شاخص‌ترین قابلیت‌های معماری و فنی این پیاده‌سازی عبارتند از:
\begin{itemize}
    \item \textbf{مسیر پیش‌رو کاملاً برداری‌شده (\lr{Fully Vectorized Forward Pass}):} پیاده‌سازی بدون هیچ‌گونه حلقه تکرار پایتونی روی توکن‌ها انجام شده و با استفاده از دستورات \lr{torch.gather} و عملگرهای \lr{scatter\_add\_} برای خبره‌های هم‌اندازه به بیشترین بازده اجرایی در کارت‌های گرافیکی دست می‌یابد.
    \item \textbf{پایداری محاسباتی در دقت‌های مختلط:} محاسبات امتیازدهی و توابع بیشینه هموار (\lr{Softmax}) در شبکه مسیریاب در فضای \lr{FP32} انجام گرفته و همگامی کامل با قالب‌های داده \lr{FP16} و \lr{BF16} حفظ می‌شود.
    \item \textbf{انطباق با نگاشت‌های دوگانه و مدرن:} پشتیبانی جامع از وزن‌های ترانهاده \lr{Conv1D} در مدل‌های خانواده \lr{GPT-2} و پروجکشن‌های ذوب‌شده و \lr{GQA} در مدل‌های نظیر \lr{Qwen/LLaMA}.
    \item \textbf{محافظت از لایه‌های حیاتی:} لایه‌های حساس شامل سر دسته‌بندی (\lr{classifier})، سر مدل زبانی (\lr{lm\_head})، ادغام‌کننده (\lr{pooler}) و لایه‌های تعبیه (\lr{embedding}) به منظور پیشگیری از فروپاشی شناختی مدل از فرآیند تبدیل به \lr{MoE} مستثنی می‌شوند.
    \item \textbf{پشتیبانی اختیاری از تطبیق‌دهنده‌های کم‌رتبه (\lr{LoRA}):} امکان تعریف پارامترهای آموزش‌پذیر \lr{LoRA} بر روی خبره‌ها و شبکه مسیریاب برای بازیابی سریع و دقیق عملکرد مدل پس از تبدیل ساختاری.
\end{itemize}
